% Auto-generated complete chapter from 28D_artstation_portfolio_benchmark.md
\chapter{28D — ArtStation / Portfolio Benchmark}
\label{complete:root_035}

\begin{sourcemeta}
\textbf{Fuente:} \href{file:///E:/Laboral/28D_artstation_portfolio_benchmark.md}{28D\_artstation\_portfolio\_benchmark.md}\\
\textbf{Seccion:} root\\
\textbf{Estado:} duplicate\\
\textbf{Rol:} Version parcial / duplicada\\
\textbf{Lineas originales:} 439\\
\textbf{Ultima modificacion:} 2026-06-11 17:52\\
\textbf{Deduplicacion conservadora:} 0 bloques repetidos omitidos; 0 caracteres repetidos no reimpresos.
\end{sourcemeta}

\section*{Resumen de orientacion}
--- tipo: modulo\_benchmark modulo: 28D\_artstation\_portfolio\_benchmark estado: generado fecha: 2026-06-11 fuente\_base: A4\_artstation\_portfolio\_benchmark perfil\_objetivo: Alexander Woodcock — Real-Time 3D / Unity WebGL / Technical Visualization ---

\section*{Contenido completo depurado}
\addcontentsline{toc}{section}{Contenido completo depurado}

\medskip\hrule\medskip


tipo: modulo\_benchmark

modulo: 28D\_artstation\_portfolio\_benchmark

estado: generado

fecha: 2026-06-11

fuente\_base: A4\_artstation\_portfolio\_benchmark

perfil\_objetivo: Alexander Woodcock — Real-Time 3D / Unity WebGL / Technical Visualization

\medskip\hrule\medskip







\subsection{1. Estado del módulo}




Este módulo convierte el intento \textbf{A4\_artstation\_portfolio\_benchmark} en criterios prácticos para el portafolio de Alexander.




La muestra recuperada es útil, pero \textbf{no es exhaustiva}. El objetivo original era 20–35 ejemplos; la extracción real dejó 12 entradas utilizables por bloqueos de ArtStation, páginas con JavaScript y portafolios que no exponían suficiente contenido técnico.




Por tanto, este módulo debe usarse como:




\begin{itemize}
\item benchmark estructural;
\item guía de presentación;
\item patrón de case studies;
\item criterio para decidir qué proyectos mostrar;
\item no como estadística representativa del mercado.
\end{itemize}




\subsection{2. Auditoría de calidad de A4}




\begin{center}
\scriptsize
\begin{longtable}{>{\raggedright\arraybackslash}p{0.305\linewidth} >{\raggedright\arraybackslash}p{0.305\linewidth} >{\raggedright\arraybackslash}p{0.305\linewidth}}
\toprule
Criterio & Resultado & Lectura \\
\midrule
\endfirsthead
\toprule
Criterio & Resultado & Lectura \\
\midrule
\endhead
Objetivo inicial & 20–35 & no cumplido \\
Ejemplos recuperados & 12 & parcial \\
ArtStation directo & débil & bloqueos \\
Blogs técnicos & fuerte & útiles \\
GitHub Pages & medio & variable \\
Portfolios visuales & medio & falta técnica \\
Evidencia usable & sí & con cautela \\
\bottomrule
\end{longtable}
\end{center}




La parte más valiosa no vino de ArtStation puro, sino de \textbf{blogs técnicos, GitHub Pages y portafolios con proyectos técnicos visibles}.




\subsection{3. Inventario consolidado}




\begin{center}
\scriptsize
\begin{longtable}{>{\raggedright\arraybackslash}p{0.184\linewidth} >{\raggedright\arraybackslash}p{0.184\linewidth} >{\raggedright\arraybackslash}p{0.184\linewidth} >{\raggedright\arraybackslash}p{0.184\linewidth} >{\raggedright\arraybackslash}p{0.184\linewidth}}
\toprule
Label & Plataforma & Rol cercano & Tipo & Fuerza \\
\midrule
\endfirsthead
\toprule
Label & Plataforma & Rol cercano & Tipo & Fuerza \\
\midrule
\endhead
Rakik & Framer & TA / Game Dev & portfolio & media \\
Potion Shader & Jettelly & Shader TA & breakdown & alta \\
Alisavakis & blog & Technical Artist & blog técnico & alta \\
Entuna & GitHub Pages & Lead TA & CV/portfolio & media \\
Rinaz & GitHub Pages & TA/Game Design & portfolio & baja \\
Múrias & GitHub Pages & Technical Artist & portfolio & media \\
Múrias Nebulas & GitHub Pages & VFX TA & project card & media \\
Múrias OpenGL & GitHub Pages & tools/rendering & project card & media \\
julhe SDF & GitHub Pages & shader dev & blog post & alta \\
julhe Optimize & GitHub Pages & shader dev & blog post & alta \\
julhe X5 & GitHub Pages & shader dev & blog post & media \\
Rakik VFX Pack & Framer & VFX TA & project card & media \\
\bottomrule
\end{longtable}
\end{center}




\subsection{4. Patrones fuertes observados}




\subsubsection{4.1. El rol debe aparecer antes que el arte}




Los perfiles útiles no se presentan solo como “3D artist”. Usan etiquetas más específicas:




\begin{mdcode}
Bloque de codigo: text
Technical Artist
Shader Developer
Pipeline & Tools Technical Artist
Game Developer & Technical Artist
VFX Technical Artist
Real-Time Look Development
\end{mdcode}




Para Alexander, el posicionamiento más defendible es:




\begin{mdcode}
Bloque de codigo: text
Real-Time 3D Developer / Unity WebGL Technical Artist
\end{mdcode}




Una variante más orientada a tesis y visualización:




\begin{mdcode}
Bloque de codigo: text
Technical Visualization Developer — Unity WebGL / Blender / Real-Time 3D
\end{mdcode}




\subsubsection{4.2. El proyecto debe abrir con resultado visual}




Los mejores ejemplos muestran primero el resultado:




\begin{itemize}
\item imagen final;
\item video corto;
\item gif;
\item comparación antes/después;
\item captura de interfaz;
\item viewport técnico;
\item breakdown visual.
\end{itemize}




Para Alexander, cada caso debería abrir con una captura o gif del resultado final antes de explicar metodología.




\subsubsection{4.3. El texto técnico debe ser específico}




Los ejemplos fuertes explican:




\begin{itemize}
\item problema técnico;
\item restricción;
\item pipeline;
\item herramientas;
\item decisiones;
\item resultado;
\item limitaciones.
\end{itemize}




Los ejemplos débiles solo muestran imágenes sin explicar qué hizo el autor.




\subsubsection{4.4. Los portfolios visuales necesitan páginas internas}




Un grid de proyectos ayuda, pero no basta. Cada card debe llevar a una página con breakdown.




Una card sin página interna no demuestra capacidad técnica. Solo demuestra estética.




\subsubsection{4.5. Los blogs técnicos funcionan como autoridad}




Los ejemplos tipo Harry Alisavakis o julhe muestran que un blog técnico puede posicionar mejor que una galería visual si explica problemas concretos.




Para Alexander, esto favorece crear páginas tipo artículo para:




\begin{itemize}
\item CAD-to-WebGL pipeline;
\item shader de corte;
\item UI técnica en Unity;
\item jerarquía de selección;
\item optimización del modelo;
\item visualización térmica heurística.
\end{itemize}




\subsection{5. Qué debe adaptar Alexander}




\subsubsection{5.1. Hero de portfolio}




Debe comunicar rol, stack y dominio en una sola pantalla.




Propuesta:




\begin{mdcode}
Bloque de codigo: text
Alexander Woodcock
Real-Time 3D Developer / Unity WebGL Technical Artist
Technical visualization, interactive 3D systems and optimized WebGL prototypes.
\end{mdcode}




Stack visible:




\begin{mdcode}
Bloque de codigo: text
Unity · WebGL · Blender · C# · UI Toolkit · Shader Graph/HLSL · 3D Optimization
\end{mdcode}




\subsubsection{5.2. Proyecto principal}




El proyecto principal debe ser \textbf{TwinSight X500}.




No debe presentarse solo como tesis. Debe presentarse como case study técnico:




\begin{mdcode}
Bloque de codigo: text
TwinSight X500 — Interactive WebGL Technical Visualization of a Drone Assembly
\end{mdcode}




Rol sugerido:




\begin{mdcode}
Bloque de codigo: text
Solo developer / technical artist / 3D pipeline owner
\end{mdcode}




Áreas demostrables:




\begin{itemize}
\item CAD-to-WebGL pipeline;
\item Blender cleanup;
\item Unity WebGL deployment;
\item part selection;
\item exploded view;
\item cross-section shader;
\item UI bottom sheet;
\item heuristic thermal view;
\item validation package.
\end{itemize}




\subsubsection{5.3. Proyectos derivados}




La tesis puede dividirse en varios mini case studies sin inventar proyectos nuevos.




\begin{center}
\scriptsize
\begin{longtable}{>{\raggedright\arraybackslash}p{0.305\linewidth} >{\raggedright\arraybackslash}p{0.305\linewidth} >{\raggedright\arraybackslash}p{0.305\linewidth}}
\toprule
Case study & Propósito & Prioridad \\
\midrule
\endfirsthead
\toprule
Case study & Propósito & Prioridad \\
\midrule
\endhead
TwinSight X500 & flagship & 1 \\
CAD-to-WebGL pipeline & pipeline técnico & 2 \\
Cross-section shader & shader/visualization & 3 \\
Exploded view system & interaction/system & 4 \\
Bottom-sheet UI & UX/tooling & 5 \\
Fastener inspection & hierarchy/detail & 6 \\
Thermal heuristic view & visual analysis & 7 \\
\bottomrule
\end{longtable}
\end{center}




\subsection{6. Estructura recomendada para cada case study}




Cada página debe seguir esta secuencia:




\begin{mdcode}
Bloque de codigo: text
1. Hero visual
2. One-line summary
3. Role / tools / scope
4. Problem
5. Constraints
6. Pipeline
7. Technical breakdown
8. Result
9. Limitations
10. Links
\end{mdcode}




\subsubsection{Plantilla breve}




\begin{mdcode}
Bloque de codigo: markdown
# Project Title

## Summary
One sentence describing the technical result.

## Role
Solo developer / technical artist / 3D pipeline owner.

## Tools
Unity, Blender, C#, WebGL, Shader Graph/HLSL, etc.

## Problem
What technical or user problem this solved.

## Constraints
Browser, mobile-first, performance, imported CAD, academic scope.

## Breakdown
Step-by-step explanation with screenshots.

## Result
What became visible, usable or measurable.

## Limitations
What the project does not claim.

## Links
Demo, GitHub, report, video.
\end{mdcode}




\subsection{7. Media mínima por proyecto}




\begin{center}
\scriptsize
\begin{longtable}{>{\raggedright\arraybackslash}p{0.305\linewidth} >{\raggedright\arraybackslash}p{0.305\linewidth} >{\raggedright\arraybackslash}p{0.305\linewidth}}
\toprule
Proyecto & Media mínima & Extra útil \\
\midrule
\endfirsthead
\toprule
Proyecto & Media mínima & Extra útil \\
\midrule
\endhead
TwinSight X500 & hero video & interaction gif \\
CAD pipeline & before/after & wireframe \\
Cross-section & gif & shader diagram \\
Exploded view & gif & hierarchy map \\
UI bottom sheet & screenshots & flow diagram \\
Thermal view & image pair & disclaimer \\
Fasteners & closeups & modular diagram \\
\bottomrule
\end{longtable}
\end{center}




\subsection{8. Profundidad técnica esperada}




\subsubsection{Nivel insuficiente}




\begin{mdcode}
Bloque de codigo: text
I made a drone viewer in Unity.
\end{mdcode}




\subsubsection{Nivel aceptable}




\begin{mdcode}
Bloque de codigo: text
I built a Unity WebGL prototype to inspect a drone assembly with part selection, exploded view and contextual technical information.
\end{mdcode}




\subsubsection{Nivel fuerte}




\begin{mdcode}
Bloque de codigo: text
I converted a CAD-derived Holybro X500 V2 assembly into an optimized Unity WebGL technical visualization prototype. The system combines normalized runtime hierarchy, part-level selection, exploded assembly reading, shader-based cross-section and bottom-sheet technical documentation for browser-based inspection.
\end{mdcode}




\subsection{9. Qué no debe copiar Alexander}




No copiar:




\begin{itemize}
\item layouts de otros portfolios;
\item capturas o assets de terceros;
\item textos de artículos de shaders;
\item etiquetas senior si no son defendibles;
\item “Technical Artist” como única etiqueta si el trabajo principal es visualización interactiva;
\item promesas no visibles en la build pública;
\item métricas sin evidencia;
\item claims de digital twin operacional si no hay telemetría real.
\end{itemize}




\subsection{10. Riesgos detectados para el portfolio de Alexander}




\subsubsection{Riesgo 1 — parecer solo académico}




Mitigación:




Presentar la tesis como producto técnico interactivo, no solo como entrega universitaria.




\subsubsection{Riesgo 2 — parecer solo modelador 3D}




Mitigación:




Separar claramente pipeline, Unity runtime, UI, shaders y WebGL.




\subsubsection{Riesgo 3 — mostrar demasiado informe y poca demo}




Mitigación:




El portfolio debe abrir con demo visual, no con texto académico.




\subsubsection{Riesgo 4 — sobrerreclamar capacidades}




Mitigación:




Usar disclaimers: visual product twin, heuristic thermal view, no FEA, no telemetry.




\subsection{11. Cambios concretos para el portfolio}




\subsubsection{Home}




Debe tener:




\begin{mdcode}
Bloque de codigo: text
- rol claro;
- hero visual del dron;
- demo link;
- GitHub link;
- 3–5 project cards;
- stack técnico;
- contacto.
\end{mdcode}




\subsubsection{Project grid}




Orden recomendado:




\begin{mdcode}
Bloque de codigo: text
1. TwinSight X500
2. CAD-to-WebGL Optimization Pipeline
3. Unity WebGL Interaction System
4. Shader-Based Cross-Section Tool
5. Technical UI / Bottom Sheet System
6. Optional: Blender Portrait / Character Study
\end{mdcode}




\subsubsection{Project pages}




Cada página debe tener:




\begin{mdcode}
Bloque de codigo: text
- 1 video/gif;
- 3–6 imágenes;
- 1 diagrama;
- 1 bloque técnico;
- 1 bloque de limitaciones;
- links a demo/repo.
\end{mdcode}




\subsection{12. Relación con GitHub}




El portfolio y GitHub deben reforzarse mutuamente.




\begin{center}
\scriptsize
\begin{longtable}{>{\raggedright\arraybackslash}p{0.305\linewidth} >{\raggedright\arraybackslash}p{0.305\linewidth} >{\raggedright\arraybackslash}p{0.305\linewidth}}
\toprule
Elemento & Portfolio & GitHub \\
\midrule
\endfirsthead
\toprule
Elemento & Portfolio & GitHub \\
\midrule
\endhead
Demo & visual primero & enlace técnico \\
README & resumen & detalle reproducible \\
Case study & explicación visual & código/manuales \\
Media & videos/gifs & screenshots \\
Evidencia & resultados & archivos fuente \\
\bottomrule
\end{longtable}
\end{center}




\subsection{13. Decisión de posicionamiento}




La etiqueta más estratégica es:




\begin{mdcode}
Bloque de codigo: text
Real-Time 3D Developer / Unity WebGL Technical Artist
\end{mdcode}




La etiqueta académica/profesional secundaria:




\begin{mdcode}
Bloque de codigo: text
Technical Visualization Developer
\end{mdcode}




La etiqueta que debe evitarse como principal:




\begin{mdcode}
Bloque de codigo: text
3D Artist
\end{mdcode}




No es falsa, pero reduce demasiado el perfil y oculta el valor técnico.




\subsection{14. Próxima tarea}




\begin{mdcode}
Bloque de codigo: text
A5_case_study_breakdown_benchmark
\end{mdcode}




Objetivo:




Recolectar y analizar ejemplos públicos de case studies técnicos con estructura fuerte para adaptar la presentación de TwinSight X500.




Prompt sugerido:




\begin{mdcode}
Bloque de codigo: text
Execute A5_case_study_breakdown_benchmark.

Objective:
Collect 15–25 public technical case studies relevant to real-time 3D, Unity WebGL, technical visualization, shader breakdowns, CAD-to-real-time pipelines, optimization, interaction systems and technical art tools.

Output structured tables only.

Fields:
- Source label
- URL
- Platform
- Project type
- Role positioning
- Opening media
- Problem statement
- Constraints shown
- Pipeline shown
- Technical breakdown depth
- Tools shown
- Metrics shown
- Limitations shown
- Links provided
- What Alexander can adapt
- What Alexander should not copy
- Confidence
- Date checked

Rules:
- Public sources only.
- Do not invent missing data.
- If unavailable, write “not found”.
- Keep table cells short.
- Include source URLs.
- Tables only.
- No strategy.
\end{mdcode}



